Executive Summary dalam sebuah capstone document berfungsi sebagai ringkasan utama yang memberikan gambaran singkat namun menyeluruh mengenai proyek akhir yang dikerjakan. Bagian ini biasanya dibaca pertama kali oleh dosen pembimbing, penguji, atau pihak luar, sehingga harus mampu menyampaikan inti dari keseluruhan laporan tanpa perlu membuka bab lain. Fokusnya adalah menjelaskan latar belakang permasalahan, tujuan proyek, serta kontribusi yang diharapkan dari capstone tersebut terhadap bidang keilmuan atau penerapan nyata.

Isi Executive Summary umumnya meliputi penjelasan singkat tentang masalah yang diangkat, metode atau pendekatan yang digunakan, serta ringkasan hasil utama yang dicapai. Jika proyek berbentuk implementasi sistem atau prototipe, maka bisa dijelaskan secara ringkas fitur utama atau inovasi yang menjadi keunggulan. Bagian ini juga dapat menyinggung keterkaitan proyek dengan kebutuhan industri atau masyarakat, sehingga nilai praktisnya lebih menonjol.

Selain itu, Executive Summary sebaiknya memuat ringkasan manfaat proyek, tantangan utama yang dihadapi, serta rekomendasi pengembangan ke depan. Dengan begitu, pembaca yang hanya membaca bagian ini sudah bisa memahami konteks, pendekatan, dan hasil dari capstone tanpa harus membaca seluruh isi dokumen. Singkatnya, Document Summary dalam capstone document harus menjadi "etalase" yang ringkas, padat, dan menarik, sekaligus mencerminkan esensi dari keseluruhan karya.